% !TEX root = ../main.tex
\section{Parámetros formales y datos discretos}
\label{ap:parametros}

Este apéndice consolida los parámetros de todos los conjuntos difusos y las muestras de las funciones de pertenencia. La intención es proveer la información necesaria para reproducir el sistema en cualquier herramienta numérica (Python con \texttt{scikit-fuzzy}, MATLAB Fuzzy Logic Toolbox, NumPy puro, etc.) sin requerir el código TypeScript.

\subsection{Tabla maestra de parámetros}

Se reproduce la tabla~\ref{tab:parametros} para conveniencia del lector:

\begin{longtable}{P{0.27\linewidth}P{0.18\linewidth}P{0.16\linewidth}P{0.27\linewidth}}
\toprule
\textbf{Variable} & \textbf{Conjunto} & \textbf{Tipo} & \textbf{Parámetros}\\
\midrule
\endfirsthead
\toprule
\textbf{Variable} & \textbf{Conjunto} & \textbf{Tipo} & \textbf{Parámetros}\\
\midrule
\endhead
\VariableParamsRows
\bottomrule
\end{longtable}

\subsection{Definición ejecutable en pseudo-Python}

\begin{lstlisting}[caption={Definicion equivalente del sistema en pseudo-Python},language=Python]
def triangular(x, a, b, c):
    if x <= a or x >= c: return 0
    if x < b: return (x - a) / (b - a)
    return (c - x) / (c - b)

def trapezoidal(x, a, b, c, d):
    if x < a or x > d: return 0
    if b <= x <= c: return 1
    if x < b: return (x - a) / (b - a) if a != b else 1
    return (d - x) / (d - c) if c != d else 1

inputs = {
  "average":       {"low": ("Z", 0,  0, 45, 60),
                    "regular": ("T", 45, 65, 80),
                    "high": ("Z", 70, 85, 100, 100)},
  "attendance":    {"low": ("Z", 0,  0, 55, 70),
                    "medium": ("T", 60, 75, 88),
                    "high": ("Z", 82, 92, 100, 100)},
  "assignments":   {"insufficient": ("Z", 0,  0, 45, 60),
                    "partial": ("T", 50, 70, 85),
                    "complete": ("Z", 78, 90, 100, 100)},
  "participation": {"low": ("Z", 0,  0, 35, 50),
                    "medium": ("T", 40, 60, 78),
                    "high": ("Z", 70, 85, 100, 100)},
  "exams":         {"deficient": ("Z", 0,  0, 45, 60),
                    "regular": ("T", 50, 68, 82),
                    "good": ("Z", 75, 88, 100, 100)},
}

output = {
  "low": ("Z", 0,  0, 20, 35),
  "medium": ("T", 25, 45, 65),
  "high": ("T", 55, 72, 88),
  "critical": ("Z", 78, 90, 100, 100),
}
\end{lstlisting}

\subsection{Muestreo numérico de las funciones de pertenencia}

Los archivos \texttt{data/membership-*.dat} contienen las muestras numéricas (paso 0.5 sobre $[0,100]$) usadas para generar las figuras~\ref{fig:mf-average} a~\ref{fig:mf-risk}. Se incluyen en el repositorio para facilitar la verificación de cada curva.
